\begin{figure}[!htbp]
\centering
\begin{tikzpicture}[font=\small,scale=1.15]
\coordinate (v1) at (0,3); \coordinate (v2) at (3,0); \coordinate (v3) at (-1,0.6);
\coordinate (q) at (0.5,1.1);
\fill[PaleGreen] (v2)--(v3)--(q)--cycle;
\fill[PaleRed] (v1)--(v3)--(q)--cycle;
\fill[PaleGrey] (v1)--(v2)--(q)--cycle;
\draw[thick,MidGrey] (v1)--(v2)--(v3)--cycle;
\foreach \v in {v1,v2,v3}{\draw[MidGrey] (q)--(\v); \fill[SapienzaRed] (\v) circle(2pt);}
\node[above] at (v1) {$V_1$}; \node[right] at (v2) {$V_2$}; \node[left] at (v3) {$V_3$};
\fill[SapienzaRed] (q) circle(2pt) node[above left] {$q$};
\node at (0.8,0.55) {$A_1$}; \node at (-0.2,1.55) {$A_2$}; \node at (1.2,1.65) {$A_3$};
\end{tikzpicture}
\caption{Triangle construction: $A_i$ is the sub-triangle opposite vertex $V_i$. Normalize by $A_1+A_2+A_3$ to obtain area weights.}
\label{fig:triangle-geometry}
\end{figure}

